\subsection{Build and Deploy: Portal}

\subsubsection{Enable Under Construction page}

Sometimes the Portal must be taken down for updates. For instance, the EC2 the Portal runs on needs to be respawned for updates.

To help facilitate communication with users, an Under Construction page can be enabled. All traffic to the Portal will be redirected to this page.

\begin{enumerate}
\item To enable the page, log into the AWS account and go to EC2 console.


\item Go to the Portal load balancer, select the \texttt{HTTPS:443} listener, and \textit{check} the Default rule.


\item In the dropdown Actions menu, select Edit Rule.


\item Set the instance target group weight to \textbf{0}. Set the lambda target group weight to \textbf{1}.


\item At the bottom of the page Save Changes.


\item Changes should take affect almost immediately.
\end{enumerate}

To revert changes after updating, repeat the above steps except change the target group weights so that the instance gets \textbf{1} and the lambda gets \textbf{0}.

\subsubsection{----------}

The following documentation is older and must be used with caution.

\subsubsection{Prerequisites}

\paragraph{AWS SES: Store SES secrets}

These secrets will be used to communicate with SES to send emails. ``SMTP credentials consist of a username and a password. When you click the Create button below, SMTP credentials will be generated for you.'' The credentials are AWS access keys, like as used in local aws configs. They are valid for the whole region. \href{https://us-west-2.console.aws.amazon.com/ses/home}{https://us-west-2.console.aws.amazon.com/ses/home}

\begin{itemize}
\item Create a verified email and take out of sandbox.


\item Create SES serets

Go to \texttt{Account Dashboard} \textgreater  \texttt{Create SMTP credentials}. The IAM User Name should be unique and easy to find within IAM. On user creation, SMTP credentials will be created.


\item Store SES secrets

\begin{enumerate}
\item Go to \href{https://us-west-2.console.aws.amazon.com/secretsmanager/home}{https://us-west-2.console.aws.amazon.com/secretsmanager/home}


\item Click on \texttt{Store New Secret} \textgreater  \texttt{Other type of secret} \textgreater  \texttt{Plaintext}


\item Delete all empty json content.


\item Add username and password as given previously in the following format: \texttt{USERNAME PASSWORD}.


\item Click \texttt{Next}


\item Secret Name: \texttt{portal/ses-creds}, Tags: \texttt{osl-billing: osl-portal}


\item Click \texttt{Next}


\item Click \texttt{Next}


\item Click \texttt{Store}
\end{enumerate}
\end{itemize}

\paragraph{AWS Secrets Manager: Create SSO token}

This token will be used by the labs to communicate and authenticate with the portal. All labs and the portal share this token. It is imperative that this remains secret. The form of the token is very specific. Use the following process to create the token.

\begin{itemize}
\item Create secret

\begin{verbatim}
pip install cryptography
\end{verbatim}

\begin{verbatim}
from cryptography.fernet import Fernet

api_token = Fernet.generate_key()
api_token
\end{verbatim}


\item Add to AWS Secret Manager

\begin{enumerate}
\item Go to \href{https://us-west-2.console.aws.amazon.com/secretsmanager/home}{https://us-west-2.console.aws.amazon.com/secretsmanager/home}


\item Click on \texttt{Store New Secret} \textgreater  \texttt{Other type of secret} \textgreater  \texttt{Plaintext}


\item Delete all empty json content.


\item Add \textit{api\_token}.


\item Click \texttt{Next}


\item Secret Name: \texttt{\$CONTAINER\_NAMESPACE/sso-token}, Tags: \texttt{osl-billing: osl-portal}


\item Click \texttt{Next}


\item Click \texttt{Next}


\item Click \texttt{Store}
\end{enumerate}
\end{itemize}

\paragraph{Docker registry}

For local dev development, one can use a local docker registry.
\texttt{docker run -d -p 5000:5000 --restart=always --name registry registry:2}

Otherwise, the remote docker images will be stored in AWS ECR, as setup by CloudFormation

\paragraph{Docker repo}

Clone the portal code: \texttt{git clone git@github.com:ASFOpenSARlab/deployment-opensciencelab-prod-portal.git}

\subsubsection{Setup}

If production, upload the Cloudformation template \texttt{cf-portal-setup.yaml} and build.

Once the cloudformation is done, go to EC2 Connect of the portal EC2, log onto the server and \texttt{cd /home/ec2-user/code}.
Then setup prerequisites via \texttt{make setup-ec2}.
Note that you will be warned about reformatting the DB volume. If this is the first time running (as it should be), do so.

If locally, go to the root of the docker repo.
Then setup prerequisites via \texttt{make setup-ubuntu}.

\subsubsection{Build}

\texttt{cp labs.example.yaml labs.\{maturity\}.yaml}. The name of the config doesn't matter (except it cannot be labs.run.yaml)
Update labs.\{maturity\}.yaml as needed

\texttt{make config=labs.\{maturity\}.yaml}

\subsubsection{Destroy}

If production, clear out the registry images, delete the CloudFormation setup, delete db ebs snapshots, and delete logs.

If locally, \texttt{make clean} and then stop the localhost registry (if being used).

\subsubsection{Other less used procedures}

\paragraph{Logs}

In production, normally the logs will show up in CloudWatch.

For both, \texttt{docker compose logs -f}.

\paragraph{Replace Portal DB from snapshot}

If the Portal DB needs to be replaced by a snapshot backup, do the following.

\textit{Unless otherwise sated, all of these steps take place within EC2 Connect.}

Elevated permissions will be needed via \texttt{sudo} or \texttt{sudo su -}.

\begin{enumerate}
\item Restore backup snapshot to volume

This procedure assumes that the usual DB volume is present and being used. We only want to update the DB file.

Within \texttt{cf-portal-setup.yaml}, it is assumed that the AZ of the EC2's subnet is set as us-west-2a.

a. From the EC2 console, select the snapshot that will be restored. Get the SNAPSHOT\_ID, e.g. snap-0c0dbee2e7c9f0c12

b. From the EC2 console, select the portal EC2. Get the EC2\_INSTANCE\_ID, e.g. i-0ca96843e97d9bd29

c. First run a dry run to make sure permissions are available for EBS volume creation.
\texttt{     aws ec2 create-volume {\textbackslash}~         --dry-run {\textbackslash}~         --availability-zone us-west-2a {\textbackslash}~         --snapshot-id \$SNAPSHOT\_ID {\textbackslash}~         --tag-specifications 'ResourceType=volume,Tags=[\{Key=Name,Value=portal-db-backup\}]'     }

\begin{verbatim}
Then actually create an EBS volume from the backup snapshot.
 ```
 aws ec2 create -volume \
     - -availability -zone us -west -2a \
     - -snapshot -id $SNAPSHOT_ID \
     - -tag -specifications 'ResourceType=volume,Tags=[{Key=Name,Value=portal -db -backup}]'
 ```

 From response output, get VOLUME_ID, e.g. vol -0a0869b5ab9f77090
\end{verbatim}


\item Attach backup volume to EC2

First run a dry run to make sure permissions are available.

\begin{verbatim}
aws ec2 attach -volume \
    - -dry -run \
    - -device /dev/sdm \
    - -instance -id $EC2_INSTANCE_ID \
    - -volume -id $VOLUME_ID
\end{verbatim}

Then actually run it.

\begin{verbatim}
aws ec2 attach -volume \
    - -device /dev/sdm \
    - -instance -id $EC2_INSTANCE_ID \
    - -volume -id $VOLUME_ID
\end{verbatim}


\item Mount device to filesystem

\texttt{sudo mkdir -p /tmp/portal-db-from-snapshot/}
\texttt{sudo mount /dev/sdm /tmp/portal-db-from-snapshot/}

If you get an error message something like

\begin{quote}
/wrong fs type, bad option, bad superblock
\end{quote}

then you cannot mount the filesystem. AWS's way of handling volumes makes things difficult. \href{https://serverfault.com/questions/948408/mount-wrong-fs-type-bad-option-bad-superblock-on-dev-xvdf1-missing-codepage}{https://serverfault.com/questions/948408/mount-wrong-fs-type-bad-option-bad-superblock-on-dev-xvdf1-missing-codepage}

Since we are working with a temporay mount, run the following instead:

\texttt{sudo mount -t xfs -o nouuid /dev/sdm /tmp/portal-db-from-snapshot/}.


\item Check for mounted directories

\texttt{df}

Look for something like (not /dev/xvdj)

\begin{verbatim}
/dev/xvdj        1038336   34224   1004112   4% /srv/jupyterhub
/dev/xvdm        1038336   34224   1004112   4% /tmp/portal -db -from -snapshot
\end{verbatim}


\item Create a backup of the old DB file within the filesystem for just in case

\texttt{sudo cp ./srv/portal/jupyterhub/jupyterhub.sqlite ./srv/portal/jupyterhub/jupyterhub.sqlite.\$(date +"\%F-\%H-\%M-\%S")}


\item Copy over backup DB file

\texttt{sudo cp /tmp/portal-db-from-snapshot/jupyterhub.sqlite ./srv/portal/jupyterhub/jupyterhub.sqlite}


\item Unmount and detach backup volume from EC2

\texttt{sudo umount /tmp/portal-db-from-snapshot/}
\texttt{sudo aws ec2 detach-volume --volume-id VOLUME\_ID}


\item Delete backup volume

\texttt{sudo aws ec2 delete-volume --volume-id VOLUME\_ID}
\end{enumerate}